\documentclass[10pt,a4paper]{amsart}
\usepackage[margin=19mm]{geometry}
\usepackage{amsmath,amssymb,mathtools}
\usepackage[scr=rsfs,cal=euler]{mathalfa}
\usepackage{xfrac}
\usepackage[unicode,hidelinks,bookmarks=false]{hyperref}
\newcommand{\ie}{i.e.\ }
\newcommand{\cf}{cf.\ }
\newcommand{\resp}{respectively\ }
\newcommand{\Subsection}{\subsection}
\providecommand{\nobreakdash}{\nobreak\hbox{-}}
\setlength{\emergencystretch}{2em}
\multlinegap0pt
\usepackage{xcolor}
\usepackage{mathtools}
\usepackage{amssymb}
\usepackage{bm}
\usepackage{csquotes}
\usepackage[shortlabels]{enumitem}
\newtheorem{lemma}{Lemma}
\usepackage{cleveref}
\crefname{lemma}{Lemma}{Lemmas}
\newcommand*\dif{\mathop{}\!\mathrm{d}}
\title{Self-normalised Bennett inequalities for Hilbert-valued martingales}
\author{David Janz}
\date{Source: arXiv:2608.15874v1 (CC BY 4.0)}
\begin{document}
\maketitle

\noindent\emph{Source and licence.} Self-normalised Bennett inequalities for Hilbert-valued martingales. Version arXiv:2608.15874v1 (CC BY 4.0), distributed by arXiv under the Creative Commons Attribution 4.0 International licence, http://creativecommons.org/licenses/by/4.0/, as recorded in the arXiv licence metadata for this exact version and shown on the abstract page https://arxiv.org/abs/2608.15874v1. Full licence notice: \texttt{licenses/arxiv-2608.15874-CC-BY-4.0.txt}.\par\smallskip

\noindent\textit{Editorial scope.} Counted unit: the article's lemma lem:chord-comparison (source0.tex lines 1367--1376). The article's complete original proof of it is delivered with it (source0.tex lines 1392--1462, one \texttt{proof} environment of 71 lines, which the article places away from the statement). No mathematics is written, repaired or re-derived: the statement and the proof are the article's own lines, in the article's own notation, and no retained line is substituted or re-flowed.  No further source-local context is needed: every cross-reference of every retained line resolves inside the two delivered spans.  Nothing was omitted from any retained span, so the package carries no omission record.  No retained line contains a citation, so the wrapper carries no bibliography and no bibliography entry.  No retained line contains a figure environment, an image-inclusion command or drawing code, so no image is delivered and the package declares no asset dependency.  Editorial formatting only, in addition: an AMS \texttt{amsart} preamble that restates the article's own declarations and macros used by the retained lines, namely the environments \texttt{lemma} on separate counters, together with the standard packages those lines need.  The retained environments print in this document as Lemma 1 in order of appearance, whereas the article prints the same items with the numbering of its own full document.  The complete native source of the article as distributed in the arXiv e-print for this exact version is bundled unchanged as \texttt{sources/207791/source0.tex}.
\begin{lemma}\label{lem:chord-comparison}
Let $0\leq\alpha\leq\beta<1$ and $r\geq0$, write $\ell=\log(1+r)$, and set
\[
  \bar r=r\frac{1-\alpha\beta}{1-\beta^2}.
\]
Then
\[
  (1+\bar r)^\beta\leq1+\beta\bar r-(r-\ell)\alpha+\frac{r\ell}{2}\alpha^2.
\]
\end{lemma}

\begin{proof}[Proof of \cref{lem:chord-comparison}]
Set $c=(1-\alpha\beta)/(1-\beta^2)$, so that $\bar r=cr$, and define
\[
  \Phi(r)=1+\beta cr-(1+cr)^\beta-(r-\ell)\alpha+\frac{r\ell}{2}\alpha^2.
\]
It suffices to show that $\Phi(r)\geq0$. We have $\Phi(0)=\Phi'(0)=0$, while
\[
  \Phi''(r)=\beta(1-\beta)c^2(1+cr)^{\beta-2}-\frac{\alpha}{(1+r)^2}+\frac{\alpha^2}{2}\left(\frac1{1+r}+\frac1{(1+r)^2}\right).
\]
Put $u=1+r$ and define
\[
  H(u):=\beta(1-\beta)c^2\frac{u^2}{\{1+c(u-1)\}^{2-\beta}}.
\]
Since
\[
  -\frac{\alpha}{u^2}+\frac{\alpha^2}{2}\left(\frac1u+\frac1{u^2}\right)\geq-\frac{\alpha(1-\alpha)}{u^2},
\]
it is enough to prove that $H(u)\geq\alpha(1-\alpha)$ for $u\geq1$.

The case $\alpha=0$ is immediate. Suppose that $0<\alpha\leq\beta<1$. The logarithmic derivative of $H$ has the sign of $\beta cu-2(c-1)$. It changes sign at most once, from negative to positive, so the minimum of $H$ over $u\geq1$ is attained either at $u=1$ or at the unique interior stationary point.

At $u=1$,
\[
  H(1)-\alpha(1-\alpha)=\frac{(\beta-\alpha)\{1-\alpha(1+\beta-\beta^2)\}}{(1-\beta)(1+\beta)^2}\geq0.
\]
Indeed, $\alpha\leq\beta$ and
\[
  1-\alpha(1+\beta-\beta^2)\geq1-\beta(1+\beta-\beta^2)=(1-\beta)^2(1+\beta)\geq0.
\]
This proves the required bound when the minimum of $H$ is attained at $u=1$.

It remains to consider the case where the stationary point of $H$ lies in $(1,\infty)$. This point is
\[
  u_\star=\frac{2(c-1)}{\beta c}=\frac{2(\beta-\alpha)}{1-\alpha\beta}.
\]
The condition $u_\star>1$ is equivalent to $2\beta-1>\alpha(2-\beta)$. Since $\alpha>0$, this can occur only when $\beta>1/2$; in that case, it is equivalent to
\[
  0<\alpha<\frac{2\beta-1}{2-\beta}=:\alpha_c.
\]
We now verify the required bound throughout this remaining range of parameters.

At the stationary point,
\[
  H(u_\star)=M_\beta(\alpha):=\frac{4\beta(1-\beta)^{1-\beta}(\beta-\alpha)^\beta}{(1+\beta)^\beta(2-\beta)^{2-\beta}}.
\]
Thus it remains to prove $M_\beta(\alpha)\geq\alpha(1-\alpha)$ for $0<\alpha<\alpha_c$. Define
\[
  Q(\alpha)=\frac{M_\beta(\alpha)}{\alpha(1-\alpha)}.
\]
Then
\[
  \frac{\dif}{\dif\alpha}\log Q(\alpha)=\frac{N_\beta(\alpha)}{\alpha(\beta-\alpha)(1-\alpha)}, \qquad N_\beta(\alpha):=(\beta-2)\alpha^2+(\beta+1)\alpha-\beta.
\]
The denominator is positive on $(0,\alpha_c]$. Moreover, $N_\beta'$ is decreasing and
\[
  N_\beta'(\alpha_c)=3(1-\beta)>0,
\]
so $N_\beta$ is increasing on $[0,\alpha_c]$. Since $N_\beta(\alpha_c)=\beta-1<0$, it follows that $N_\beta(\alpha)<0$ throughout this interval. Hence $Q$ is decreasing on $(0,\alpha_c]$.

Finally,
\[
  M_\beta(\alpha_c)=\frac{4\beta(1-\beta)}{(2-\beta)^2}, \qquad \alpha_c(1-\alpha_c)=\frac{3(2\beta-1)(1-\beta)}{(2-\beta)^2},
\]
and therefore
\[
  Q(\alpha_c)=\frac{4\beta}{3(2\beta-1)}\geq1.
\]
Since $Q$ is decreasing, $Q(\alpha)\geq Q(\alpha_c)\geq1$ for every $0<\alpha<\alpha_c$. Thus $H(u_\star)\geq\alpha(1-\alpha)$.

We have now controlled the minimum of $H$ whether it is attained at $u=1$ or at the interior stationary point. Hence $\Phi''(r)\geq0$ for every $r\geq0$. Since $\Phi(0)=\Phi'(0)=0$, it follows that $\Phi(r)\geq0$.
\end{proof}

\end{document}
